\section{Гранични теореми и приближения}

\subsection{Теорема на Поасон (Поасоново приближение)}

Досега разглеждахме Поасоновата случайна величина като модел, дефиниран чрез своите вероятности. За разлика от биномното разпределение, при което имаме ясен експериментален модел (хвърляне на монета или повтаряеми опити с два изхода), Поасоновото разпределение често възниква като граничен случай или като добро приближение на други разпределения.

Преди да формулираме основната теорема, нека припомним едно много полезно твърдение, свързано с пораждащите функции. Нека имаме редица от целочислени неотрицателни случайни величини $X_n$ и съответните им пораждащи функции $G_{X_n}(s)$. Нека $X$ е също целочислена неотрицателна случайна величина с пораждаща функция $G_X(s)$. Ако за всяко $s$ от дефиниционното множество имаме сходимост на пораждащите функции $\lim_{n \to \infty} G_{X_n}(s) = G_X(s)$, то за всяко цяло $k \ge 0$ е изпълнено и сходимост на самите вероятности:
\[
\lim_{n \to \infty} \mathbb{P}(X_n = k) = \mathbb{P}(X = k)
\]
Това твърдение ни позволява чрез граница на функции да идентифицираме граница на вероятности, което много често е технически значително по-лесно от директното доказване на сходимост на самите вероятности.

Ще използваме този факт, за да докажем една от важните ранни теореми в теорията на вероятностите, а именно Теоремата на Поасон. Тя илюстрира как един теоретичен резултат може да ни даде основание за полезни практически приближения.

\begin{keythm}[на Поасон]
Нека за всяко $n \ge 1$ е дадена случайна величина $X_n$, която има биномно разпределение $X_n \sim \Bin(n, p_n)$. Т.е. имаме редица от биномни случайни величини с нарастващ брой експерименти $n$ и променяща се вероятност за успех $p_n$. Нека освен това границата на произведението $n p_n$ е положително число:
\[
\lim_{n \to \infty} n p_n = \lambda > 0
\]
Тогава за всяко $k \ge 0$, границата на вероятностите $X_n$ да е равно на $k$ е равна на вероятността Поасонова случайна величина с параметър $\lambda$ да е равна на $k$:
\[
\lim_{n \to \infty} \mathbb{P}(X_n = k) = \frac{\lambda^k}{k!} e^{-\lambda}
\]
Грубо казано, ако броят експерименти по вероятността за успех се стабилизира до някакво число $\lambda$, сложните биномни вероятности се схождат към много по-лесно изчислимите Поасонови вероятности.
\end{keythm}

\begin{proof}[Доказателство (за частен случай)]
Няма да доказваме теоремата в пълната ѝ общност. За да предадем основната интуиция, ще разгледаме по-простия случай, при който не просто границата е $\lambda$, а самото произведение е точно равно на $\lambda$ за всяко $n$. Т.е. нека $\lambda = n p_n$, откъдето $p_n = \frac{\lambda}{n}$.

Пораждащата функция на биномната случайна величина $X_n$ е:
\[
G_{X_n}(s) = \left(1 - p_n + p_n s\right)^n
\]
Замествайки $p_n = \frac{\lambda}{n}$, получаваме:
\[
G_{X_n}(s) = \left(1 - \frac{\lambda}{n} + \frac{\lambda}{n} s\right)^n = \left(1 + \frac{\lambda (s-1)}{n}\right)^n
\]
Като използваме основното свойство на експоненциалната функция като граница, намираме:
\[
\lim_{n \to \infty} G_{X_n}(s) = \lim_{n \to \infty} \left(1 + \frac{\lambda (s-1)}{n}\right)^n = e^{\lambda (s-1)}
\]
Получената граница $e^{\lambda (s-1)}$ е точно пораждащата функция на Поасонова случайна величина $X \sim \Poi(\lambda)$. Благодарение на споменатото по-горе твърдение, това доказва, че $\lim_{n \to \infty} \mathbb{P}(X_n = k) = \frac{\lambda^k}{k!} e^{-\lambda}$.

\textbf{Практическо приложение (Приближение на Поасон):}
В практиката тази теорема не се използва толкова като строга граница (понеже на практика разполагаме с едно конкретно $n$), колкото като рецепта за приближение. Сложният израз за биномната вероятност $\binom{n}{k} p^k (1-p)^{n-k}$ изисква пресмятането на големи факториели и високи степени.

Ако разполагаме със случайна величина $Y \sim \Bin(n, p)$, където $n$ е достатъчно голямо, а $p$ е достатъчно малко, можем да използваме приближението:
\[
\mathbb{P}(Y = k) \approx \frac{\lambda^k}{k!} e^{-\lambda}, \quad \text{където } \lambda = n p
\]
\emph{Рецепта:} Приближението работи добре и може спокойно да се използва в случаите, когато $n \ge 100$ и $np \le 20$. Тогава, дори и да не знаем „границите“, твърдим, че биномните вероятности са изключително близки до Поасоновите с параметър $\lambda = np$.
\end{proof}


\section{Хипергеометрично разпределение}

Последното конкретно разпределение от фундаменталните дискретни случайни величини, което ще разгледаме, е хипергеометричното разпределение. За разлика от Поасоновото, то е директно обвързано с един много ясен физически и комбинаторен модел на извличане без връщане.

\begin{defn}[хипергеометричен модел]
Казваме, че $X$ има хипергеометрично разпределение с параметри $N$, $M$ и $n$ (записваме $X \sim \Hyp(N, M, n)$), ако $X$ описва следния процес:
Имаме общо $N$ на брой обекта (например продукти в партида). От тях точно $M$ са „маркирани“ (или „дефектни“), а останалите $N - M$ са немаркирани (или здрави). Ние изтегляме случайно $n$ на брой обекта от тези $N$ без връщане. Случайната величина $X$ представлява \emph{броят на изтеглените маркирани обекти} измежду тези $n$ обекта.
\end{defn}

Естествените параметрични ограничения са $N \ge M$ и $N \ge n$. Това разпределение намира огромно приложение в качествения контрол, където искаме да оценим каква пропорция от огромна партида стоки са дефектни, използвайки само малка извадка от тях.

\begin{prop}
Нека $X \sim \Hyp(N, M, n)$. Тогава:
\begin{enumerate}
    \item[а)] Вероятностното разпределение на $X$ се задава с:
    \[
    \mathbb{P}(X = k) = \frac{\binom{M}{k} \binom{N-M}{n-k}}{\binom{N}{n}}
    \]
    където $k$ удовлетворява естествените ограничения $0 \lor (n - N + M) \le k \le M \land n$.
    \item[б)] Математическото очакване е:
    \[
    \E X = n \frac{M}{N}
    \]
    \item[в)] Дисперсията е:
    \[
    \Var X = n \frac{M}{N} \frac{N-M}{N} \frac{N-n}{N-1}
    \]
\end{enumerate}
\end{prop}

\begin{proof}
\emph{За а):} Формулата следва от класическата вероятност. Всички възможни начини да изтеглим $n$ обекта от $N$ са $\binom{N}{n}$. За да изтеглим точно $k$ маркирани, трябва да изберем $k$ обекта от наличните $M$ маркирани (това става по $\binom{M}{k}$ начина) и останалите $n-k$ обекта да изберем от наличните $N-M$ немаркирани (това става по $\binom{N-M}{n-k}$ начина). Границите за $k$ гарантират, че не избираме повече маркирани, отколкото съществуват или отколкото теглим ($k \le \min(M, n)$), както и че няма как да изтеглим отрицателен брой обекти или повече немаркирани, отколкото съществуват.

\emph{За б):} Това свойство има скрита линейна структура, която изключително улеснява пресмятането на очакването. Дефинираме индикаторни случайни величини $X_j$ за всяко $j \in \{1, \dots, n\}$ по следния начин:
\[
X_j = \begin{cases}
1, & \text{ако } j\text{-тият изтеглен обект е маркиран} \\
0, & \text{ако не е маркиран}
\end{cases}
\]
Ясно е, че общият брой маркирани обекти $X$ е просто сумата на тези индикатори:
\[
X = \sum_{j=1}^n X_j
\]
Тъй като математическото очакване винаги е линейно (независимо дали събираемите са независими или не), можем да запишем:
\[
\E X = \E\left(\sum_{j=1}^n X_j\right) = \sum_{j=1}^n \E X_j
\]
Всяко $X_j$ е Бернулиева случайна величина $X_j \sim \Ber(p_j)$, където $p_j$ е вероятността $j$-тият изтеглен обект да е маркиран. Поради пълната симетрия на модела (няма значение дали теглите първи или $j$-ти, преди да сте видели резултатите от останалите тегления, вероятността е една и съща), вероятността на коя да е позиция да се падне маркиран обект е точно пропорцията на маркираните обекти в началото, а именно $p_j = \frac{M}{N}$.
Следователно, $\E X_j = p_j = \frac{M}{N}$. Замествайки обратно в сумата, получаваме:
\[
\E X = \sum_{j=1}^n \frac{M}{N} = n \frac{M}{N}
\]
\emph{За в):} Тук вече не можем да съберем дисперсиите наум, защото случайните величини $X_j$ \emph{не са независими} (изтеглянето на маркиран обект на първо теглене променя вероятностите за следващите тегления, тъй като тегленето е без връщане). Затова дисперсията на сумата не е просто сума от дисперсиите, а включва и ковариационните членове. Формулата се дава наготово.

\end{proof}

\begin{example}[Цветни маркери]\label{ex:07-1}
Имаме 8 маркера, от които 4 са цветни и 4 са черни. Т.е. $N = 8, M = 4$. Теглим $n = 3$ маркера на случаен принцип. Броят цветни маркери, които сме изтеглили сред тези три, се описва с хипергеометричната случайна величина $X \sim \Hyp(8, 4, 3)$.
\end{example}


\section{Съвместни разпределения}

Досега разглеждахме случайните величини изолирано. Често пъти обаче искаме да описваме сложни модели с повече от една случайна характеристика. Например, искаме да изследваме както разхода на един клиент в магазина, така и неговия пол. Затова трябва да въведем понятия за съвместна работа с две или повече случайни величини. За удобство на записа ще разглеждаме двойки $(X, Y)$, но концепциите се обобщават за произволен краен брой.

\subsection{Съвместни дискретни разпределения}

\begin{defn}
Нека $X$ и $Y$ са две дискретни случайни величини, дефинирани върху едно и също вероятностно пространство. Под \textbf{съвместно разпределение} на $X$ и $Y$ разбираме съвкупността от всички вероятности:
\[
p_{i,j} = \mathbb{P}(X = x_i, Y = y_j)
\]
за всички възможни стойности $x_i$ на $X$ и $y_j$ на $Y$. Обикновено това се представя във формата на таблица, където редовете отговарят на възможните стойности на $X$, а колоните — на възможните стойности на $Y$. Всяка клетка $(i, j)$ съдържа вероятността двете събития да настъпят едновременно. Сумата на всички числа (клетки) в тази таблица е равна на 1.
\end{defn}

\textbf{Маргинални разпределения:}
Ако съберем всички вероятности по даден ред (фиксирайки $X = x_i$ и сумирайки по всички възможни стойности на $Y$), по формулата за пълната вероятност получаваме \emph{маргиналното разпределение} на $X$:
\[
\mathbb{P}(X = x_i) = \sum_{j} \mathbb{P}(X = x_i, Y = y_j) = \sum_{j} p_{i,j}
\]
Това е просто индивидуалното разпределение на $X$, такова каквото би било, ако въобще не се интересувахме от $Y$. Аналогично, сумирайки по колони, получаваме маргиналното разпределение на $Y$:
\[
\mathbb{P}(Y = y_j) = \sum_{i} \mathbb{P}(X = x_i, Y = y_j) = \sum_{i} p_{i,j}
\]

\begin{example}\label{ex:07-2}
Хвърляме два стандартни зара (номерирани от 1 до 6). Дефинираме две случайни величини:
\begin{itemize}
    \item $X$: броят на хвърлените шестици (възможни стойности: 0, 1, 2);
    \item $Y$: броят на хвърлените единици (възможни стойности: 0, 1, 2).
\end{itemize}

Нека попълним таблицата на съвместното им разпределение. Общият брой възможни изходи при хвърляне на два зара е $6 \times 6 = 36$.
\begin{itemize}
    \item $\mathbb{P}(X=0, Y=0)$: Търсим вероятността да няма нито една шестица и нито една единица. Значи и двата зара трябва да покажат числа между 2 и 5 (общо 4 възможности за всеки зар). Вероятността е $\frac{4 \times 4}{36} = \frac{16}{36}$.
    \item $\mathbb{P}(X=1, Y=0)$: Търсим една шестица и нула единици. За първия зар имаме шестица, а за втория число от 2 до 5 (4 възможности). Или обратното. Общо $2 \times 4 = 8$ възможности. Вероятността е $\frac{8}{36}$.
    \item $\mathbb{P}(X=0, Y=1)$: По симетрия (една единица и нула шестици), вероятността също е $\frac{8}{36}$.
    \item $\mathbb{P}(X=1, Y=1)$: Една шестица и една единица. Вариантите са (6,1) и (1,6). Вероятността е $\frac{2}{36}$.
    \item $\mathbb{P}(X=2, Y=0)$: Две шестици, вероятност $\frac{1}{36}$.
    \item $\mathbb{P}(X=0, Y=2)$: Две единици, вероятност $\frac{1}{36}$.
    \item Всички останали комбинации (напр. $X=2, Y=1$) са невъзможни (не може да имаме общо повече от 2 зара), така че вероятността им е 0.
\end{itemize}

Сумата от всички вероятности в таблицата е $\frac{16+8+8+2+1+1}{36} = \frac{36}{36} = 1$.

Ако изчислим маргиналното разпределение на $X$ чрез сумиране по редовете, получаваме:
\[
\mathbb{P}(X=0) = \frac{16+8+1}{36} = \frac{25}{36}
\]
\[
\mathbb{P}(X=1) = \frac{8+2}{36} = \frac{10}{36}
\]
\[
\mathbb{P}(X=2) = \frac{1}{36}
\]
Което е точно очакваното биномно разпределение за броя на успехите (падане на шестица, $p = 1/6$) при 2 опита.


\end{example}

\subsection{Съвместна функция на разпределение (Joint CDF)}

\begin{defn}
Нека $X$ и $Y$ са две произволни случайни величини (дискретни или непрекъснати). Тогава \textbf{съвместна функция на разпределение} (или кумулативна функция) на $X$ и $Y$ е функцията $F_{X,Y} : \mathbb{R}^2 \to [0, 1]$, зададена като:
\[
F_{X,Y}(x,y) := \mathbb{P}(X < x, Y < y) \quad \forall (x,y) \in \mathbb{R}^2
\]
Геометрично, ако разглеждаме стойностите на $X$ и $Y$ като координати на двумерен вектор в равнината, $F_{X,Y}(x,y)$ е вероятността този случаен вектор да попадне в безкрайния правоъгълник на югозапад от точката $(x, y)$ (строго по-малко от $x$ по хоризонтала и строго по-малко от $y$ по вертикала).
\end{defn}

\begin{prop}[Маргинални функции от съвместната]\label{prop:marginal-from-joint}
От съвместната функция на разпределение се възстановяват маргиналните:
\[
F_{X,Y}(x, \infty) = F_X(x), \qquad F_{X,Y}(\infty, y) = F_Y(y), \qquad F_{X,Y}(\infty, \infty) = 1 .
\]
\end{prop}

Възстановяването става, като заместим съответно другата координата с безкрайност. Понеже събитието $\{Y < \infty\}$ е достоверно (вероятност 1), то сечението на $\{X < x\}$ с достоверно събитие не променя вероятността:
\[
F_{X,Y}(x, \infty) = \mathbb{P}(X < x, Y < \infty) = \mathbb{P}(X < x) = F_X(x)
\]
Аналогично:
\[
F_{X,Y}(\infty, y) = F_Y(y)
\]
Очевидно е също, че $F_{X,Y}(\infty, \infty) = \mathbb{P}(X < \infty, Y < \infty) = 1$.

\begin{keythm}[Критерий за независимост]\label{thm:indep-criterion}
Две случайни величини $X$ и $Y$ са независими тогава и само тогава, когато съвместната им функция на разпределение се разпада на произведение от маргиналните им функции на разпределение за всяка точка $(x,y)$:
\[
F_{X,Y}(x, y) = F_X(x) F_Y(y) \quad \forall (x,y) \in \mathbb{R}^2 .
\]
\end{keythm}

Това е най-общият критерий за независимост. За дискретни разпределения обикновено се ползва еквивалентният и по-интуитивен критерий с вероятности в конкретни точки: $\mathbb{P}(X = x, Y = y) = \mathbb{P}(X = x) \mathbb{P}(Y = y)$.

\begin{remark}[Сума на независими величини и зависимост]\label{rem:sum-dependence}
Ако случайните величини $X$ и $Y$ са независими и стандартно нормално разпределени ($N(0,1)$), а $Z$ зависи от $Y$, следва ли, че $X+Y$ е независимо от $Z$?

Не — и контрапримерът е елементарен. Нека $Y = Z$. Тогава сумата $X+Y$ е всъщност $X+Z$, което очевидно зависи директно от $Z$. Може да има специфични случаи на независимост, но в общия случай $X+Y$ и $Z$ са зависими.
\end{remark}


\section{Ковариация и корелация}

След като дефинирахме какво означава две случайни величини да са независими във вероятностен смисъл, преминаваме към изследването на конкретни видове зависимости. Понякога знаем, че две величини са зависими, но искаме да измерим силата и посоката на тяхната \emph{линейна} зависимост. Търсим да отговорим на въпроса доколко $Y$ може да се представи като линейна функция на $X$, т.е. $Y \approx aX + b$.

\subsection{Ковариация}

Първата стъпка към дефиниране на такава мярка е ковариацията. За да изолираме отместванията на разпределенията по реалната права, ние първо ги центрираме, като изваждаме техните математически очаквания.
Ако дефинираме центрирани случайни величини $\widetilde{X} = X - \E X$ и $\widetilde{Y} = Y - \E Y$, то е очевидно, че тяхното очакване е 0 ($\E\widetilde{X} = 0$).

\begin{defn}
Ковариация на $X$ и $Y$ се нарича математическото очакване на произведението на техните центрирани версии (ако това очакване съществува):
\[
\Cov(X, Y) = \E[(X - \E X)(Y - \E Y)] = \E[\widetilde{X}\widetilde{Y}]
\]
Ковариацията може да се изчислява и по алтернативна формула, която се извежда директно чрез разкриване на скобите и използване на линейността на математическото очакване:
\[
\Cov(X, Y) = \E[XY] - \E X \E Y
\]
\end{defn}
\begin{cor}
Ако $X$ и $Y$ са независими (най-силното свойство на липса на зависимост), то математическото очакване на произведението им се разпада на произведение от очакванията им ($\E[XY] = \E X\E Y$). Следователно, при независими случайни величини ковариацията е точно 0. Важно е да се отбележи, че обратното не е вярно — нулева ковариация означава само липса на линейна зависимост, но величините може да са свързани чрез по-сложна, нелинейна функция.
\end{cor}

\begin{prop}[Дисперсия на сума]\label{prop:var-sum}
За произволни случайни величини с крайни дисперсии
\[
\Var(aX + bY) = a^2 \Var X + b^2 \Var Y + 2ab\,\Cov(X,Y),
\]
и по-общо
\[
\Var\Big( \sum_{j=1}^{n} X_j \Big)
 = \sum_{j=1}^{n} \Var X_j + 2 \sum_{i<j} \Cov(X_i, X_j) .
\]
\end{prop}

\begin{proof}
Достатъчно е да разкрием квадрата под очакването. Полагаме $\tilde X = X - \E X$ и
$\tilde Y = Y - \E Y$; тогава $a X + b Y$ има центрирана част $a\tilde X + b\tilde Y$ и
\[
\Var(aX+bY) = \E\big[(a\tilde X + b\tilde Y)^2\big]
 = a^2 \E \tilde X^2 + b^2 \E \tilde Y^2 + 2ab\,\E[\tilde X \tilde Y],
\]
което е точно търсеното. Общият случай се получава по същия начин: при повдигане на
$\sum_j \tilde X_j$ на квадрат диагоналните членове дават дисперсиите, а всяка
недиагонална двойка се появява два пъти и дава $2\Cov(X_i, X_j)$.
\end{proof}

Ако величините са независими в съвкупност, всички ковариации се нулират и се връщаме към
адитивността на дисперсията от лекция 5. Точно тези „ковариационни членове“ липсват при
хипергеометричното разпределение, където изтеглянията не са независими.

\emph{Недостатък на ковариацията:} Ако умножим случайните величини по константа, например 10 пъти, тяхната ковариация ще се увеличи 100 пъти: $\Cov(10X, 10Y) = 100 \Cov(X, Y)$. Това математическо мащабиране променя стойността на мярката, въпреки че линейната структура на зависимостта между величините реално не се е променила.

\subsection{Коефициент на корелация}

За да компенсираме ефекта от мащабирането, ние не само центрираме случайните величини, но ги и \emph{нормираме}, разделяйки ги на корен квадратен от техните дисперсии (т.е. на стандартните им отклонения).

\begin{defn}
Коефициент на корелация $\rho(X, Y)$ между две случайни величини с крайни дисперсии $\Var X$ и $\Var Y$ се дефинира като:
\[
\rho(X, Y) = \frac{\Cov(X, Y)}{\sqrt{\Var X} \sqrt{\Var Y}}
\]
Ако дефинираме центрираните и нормирани случайни величини като:
\[
\overline{X} = \frac{X - \E X}{\sqrt{\Var X}}, \quad \overline{Y} = \frac{Y - \E Y}{\sqrt{\Var Y}}
\]
може лесно да се провери, че $\E\overline{X} = 0$ и $\Var\overline{X} = 1$. В тези термини, коефициентът на корелация е просто очакването на тяхното произведение:
\[
\rho(X, Y) = \E[\overline{X}\overline{Y}]
\]
Коефициентът на корелация е безразмерна величина. Ако пресметнем корелацията на $(10X, 10Y)$, множителят $100$ се появява както в числителя (ковариацията), така и в знаменателя (от корените на дисперсиите), и се съкращава. Следователно мащабирането не променя корелацията.
\end{defn}

\begin{keythm}[Граници на корелацията]\label{thm:corr-bounds}
За произволни случайни величини с крайни дисперсии е вярно, че:
\[
|\rho(X, Y)| \le 1 .
\]
\end{keythm}

Това е вероятностната форма на \emph{неравенството на Коши\,--\,Буняковски} (в
литературата и като неравенство на Коши\,--\,Шварц). Връзката с геометрията е
пряка: ковариацията играе ролята на скаларно произведение, стандартните
отклонения — на дължини, а коефициентът на корелация — на косинус от ъгъла между
двата вектора, откъдето и границите $\pm 1$.
\begin{proof}
Ще използваме очевидния факт, че квадратът на всяка случайна величина е неотрицателен, следователно и очакването му е неотрицателно:
\[
0 \le \E[(\overline{X} \pm \overline{Y})^2]
\]
Разкриваме скобите, ползвайки линейността на очакването:
\[
0 \le \E[\overline{X}^2 \pm 2\overline{X}\overline{Y} + \overline{Y}^2] = \E\overline{X}^2 \pm 2\E[\overline{X}\overline{Y}] + \E\overline{Y}^2
\]
Тъй като $\overline{X}$ и $\overline{Y}$ имат нулеви очаквания, техните дисперсии съвпадат с очакването на квадратите им: $\E\overline{X}^2 = \Var\overline{X} = 1$ и $\E\overline{Y}^2 = \Var\overline{Y} = 1$. Заместваме това заедно с дефиницията за $\rho(X, Y) = \E[\overline{X}\overline{Y}]$ и получаваме:
\[
0 \le 1 \pm 2\rho(X, Y) + 1 = 2 \pm 2\rho(X, Y)
\]
Следователно $1 \pm \rho(X, Y) \ge 0$, откъдето правим извода, че $|\rho(X, Y)| \le 1$.
\end{proof}

Следващата теорема обяснява точно какво ни казва коефициентът на корелация и кога той достига своите екстремални стойности $\pm 1$.

\begin{keythm}[Критерий за линейна зависимост]
Нека $X$ и $Y$ са случайни величини с крайни дисперсии. Тогава между тях има строго линейна зависимост от вида $Y = aX + b$ (за някакви константи $a \neq 0$ и $b$) тогава и само тогава, когато $|\rho(X, Y)| = 1$.
\end{keythm}

\begin{proof}[Доказателство ($\implies$)]
Нека е дадено, че $Y = aX + b$. Ще изразим връзката между нормираните величини $\overline{Y}$ и $\overline{X}$.
\[
\overline{Y} = \frac{Y - \E Y}{\sqrt{\Var Y}} = \frac{aX + b - (a\E X + b)}{\sqrt{\Var Y}} = \frac{a}{\sqrt{\Var Y}}(X - \E X)
\]
Тъй като $\sqrt{\Var X} = \frac{\sqrt{\Var Y}}{|a|}$, следва че:
\[
\overline{Y} = \frac{a}{|a|} \frac{X - \E X}{\sqrt{\Var X}} = v \overline{X}
\]
където $v = \text{sgn}(a) = \pm 1$. Така получихме връзката $\overline{Y} = \pm \overline{X}$.
Сега пресмятаме коефициента на корелация:
\[
\rho(X, Y) = \E[\overline{X}\overline{Y}] = \E[\overline{X}(\pm \overline{X})] = \pm \E\overline{X}^2 = \pm 1
\]
Тъй като $\E\overline{X}^2 = \Var\overline{X} = 1$. Това показва, че $|\rho(X, Y)| = 1$.
\end{proof}

\begin{proof}[Доказателство ($\Longleftarrow$)]
Нека сега е дадено, че $|\rho(X, Y)| = 1$. Ще разгледаме случая $\rho(X, Y) = 1$ (доказателството за $\rho = -1$ е аналогично).
От предишното доказателство знаем, че:
\[
\E[(\overline{X} - \overline{Y})^2] = 2 - 2\rho(X, Y)
\]
При $\rho(X, Y) = 1$, получаваме $\E[(\overline{X} - \overline{Y})^2] = 0$.
Тъй като интегрираме (търсим математическо очакване) на строго неотрицателна случайна величина и получаваме 0, това означава, че самата случайна величина е тъждествено равна на нула (с вероятност 1):
\[
(\overline{X} - \overline{Y})^2 = 0 \implies \overline{X} = \overline{Y}
\]
Заместваме дефинициите за центрирани и нормирани величини:
\[
\frac{X - \E X}{\sqrt{\Var X}} = \frac{Y - \E Y}{\sqrt{\Var Y}}
\]
От това уравнение лесно можем да изразим $Y$ като линейна функция на $X$:
\[
Y = \frac{\sqrt{\Var Y}}{\sqrt{\Var X}} X - \frac{\sqrt{\Var Y}}{\sqrt{\Var X}} \E X + \E Y
\]
Което е точно от вида $Y = aX + b$, където $a = \frac{\sqrt{\Var Y}}{\sqrt{\Var X}}$ и $b = \E Y - a\E X$. С това теоремата е доказана.
\end{proof}

В заключение, коефициент на корелация близо до 1 (или -1) индикира силна права (или обратна) линейна зависимост между величините. Коефициент близо до 0 означава липса на линейна зависимост.

\section{Задачи}

\begin{enumerate}
    \item За фиксирано $k$ докажете директно границата от теоремата на Поасон:
    \[
    \lim_{n\to\infty}\binom{n}{k}\left(\frac{\lambda}{n}\right)^k
    \left(1-\frac{\lambda}{n}\right)^{n-k}
    =e^{-\lambda}\frac{\lambda^k}{k!}.
    \]
    \item За $X\sim\Hyp(N,M,n)$ докажете оставената като техническо упражнение формула
    \[
    \Var X=n\frac{M}{N}\frac{N-M}{N}\frac{N-n}{N-1}.
    \]
    \item В примера с осемте маркера, четири от които са цветни, се теглят три маркера без връщане. Пресметнете вероятностите да бъдат изтеглени съответно $0$, $1$, $2$ и $3$ цветни маркера.
    \item От дефиницията на съвместната функция на разпределение проверете
    \[
    F_{X,Y}(x,\infty)=F_X(x),\qquad
    F_{X,Y}(\infty,y)=F_Y(y).
    \]
    \item Разкрийте скобите в дефиницията на ковариацията и докажете
    \[
    \Cov(X,Y)=\E[XY]-\E X\,\E Y.
    \]
    \item[Домашно 2.] Нека $X$ и $Y$ са независими стандартно нормално разпределени случайни величини, а $Z$ зависи от $Y$. Следва ли, че $X+Y$ е независимо от $Z$? Разгледайте като контрапример случая $Y=Z$ и разпишете зависимостта. Решението е обсъдено в забележката към раздела за независимост.
\end{enumerate}
